\chapter{Introduction}
\label{ch:intro}

Questionnaire surveys are a primary source of evidence in the social sciences and
in public policy, but they are expensive, slow, and difficult to field at scale.
Large language models (LLMs) have been proposed as virtual survey respondents:
given a demographic profile and a question, a model can produce an answer without
recruiting a participant \citep{argyle2022,dillion2023,kim2023}. Most existing
work evaluates the model as a one-shot prompt. A profile is placed in the context
window, a question is asked, and the answer is scored. This measures how well a
model can role-play a respondent in a single turn. It does not test whether the
model has internalized the structure of survey responses, whether its outputs are
coherent at the population level, or whether the behaviour survives a change in
the political environment.

This thesis takes a data-centric view. Survey responses are generated by a
structured process: partisan identity, ideological position, and demographics
interact with the attributes of a candidate, and the resulting evaluations follow
regularities that are stable within an era but can shift across eras. If an LLM
is to act as a credible respondent, it should be trained on that structure rather
than prompted with it. The thesis reconstructs voter--candidate interactions from
the American National Election Studies (ANES), aligns a compact 8B model with the
resulting distribution through LoRA fine-tuning, and evaluates the model at the
individual, population, and temporal levels.

The thesis makes four contributions. First, it provides a reproducible data
construction and audit. The raw ANES table contains duplicate thermometer items
stored under trailing-period names (\texttt{obama.} and \texttt{obama}) and
social-group thermometers stored in the same column as candidate thermometers.
Both defects are quantified and removed, reducing 132,948 raw rows to 98,159
clean interaction-level rows. Because ANES is a panel (3,080 of 3,165 respondents
appear in more than one wave, and there are no test-only respondents in
2022/2024), two complementary protocols are reported: a temporal split and a
respondent-disjoint split.

Second, the thesis reports a conservative individual-level evaluation. On unseen
respondents in 2016--2019, the fine-tuned model reaches an MAE of 20.84, a
correlation of 0.679, and a choice accuracy of 0.518 against chance 0.163. A
gradient-boosted baseline on the same features reaches 20.58, 0.714, and 0.534;
paired bootstrap shows no significant MAE or choice difference, while the
baseline has a higher correlation. The value of the LLM is not superior point
prediction; it is a unified interface that also produces a rating distribution
and extends to population-level generation.

Third, on GSS policy attitudes, fine-tuning reduces marginal and joint
Jensen--Shannon divergence by 91\% and 86\% on the EDC configuration, and by up
to 76\% and 64\% on EHC. The gain comes from supervised fine-tuning, not from
retrieval or quota sampling.

Fourth, the thesis reports a temporal boundary. Across 2016--2024,
individual-level correlation falls from 0.722 in 2019 to $-0.056$ in 2024, and
the tabular baseline collapses in the same way. Adding the 2022 wave to training
does not help. Structured alignment is necessary but not sufficient for social
simulation.

\section{Related Work}

Work on LLMs as virtual respondents conditions models on demographic backstories
and compares the resulting ``silicon samples'' with survey margins
\citep{argyle2022,santurkar2023,dillion2023,kim2023}. Later work extends the
paradigm to random silicon sampling \citep{sun2024}, cross-survey probing
\citep{zhou2024}, multi-domain benchmarks \citep{zhao2025}, and specialization to
survey response distributions \citep{cao2025}. Most of this work evaluates a
model as a one-shot respondent and rarely separates individual accuracy from
population-level alignment.

A related literature studies persona fidelity and personalization
\citep{park2024,wang2024,xu2024,salemi2024,ng2024}. That work is primarily about
behavioural consistency within a conversation; the present setting asks the model
to reproduce the statistical structure of survey responses across a population.
Fine-tuning has also been used to align models with political or social
preferences \citep{chu2023,zhao2023}; here the goal is competence alignment---the
rules that map a profile and a candidate to a rating---rather than adopting a
group's voice. For population-level generation, iterative proportional fitting
\citep{deming1940,fienberg1970} and LOESS \citep{cleveland1979} provide the
classical tools that are combined with the language model.

Prior work on LLM survey simulation has generally evaluated on a single wave or on
a pooled sample, which can hide temporal decay. The per-year analysis in
Chapter~\ref{ch:temporal} reports a monotonic decay for both a fine-tuned LLM and
a feature-engineered baseline on the same ANES panel. This is, to our knowledge,
the first such comparison on a common panel.
